%! AP Statistics — Unit 1, Lesson 4
\documentclass[10pt]{article}
\usepackage[margin=0.6in]{geometry}
\usepackage{xcolor}
\usepackage[most]{tcolorbox}
\usepackage{enumitem}
\usepackage{multicol}
\usepackage{amsmath,amssymb}
\usepackage{array}
\usepackage{tabularx}
\tcbuselibrary{breakable}

\definecolor{sky}{HTML}{EAF3FB}
\definecolor{navy}{HTML}{1F3A5F}
\definecolor{redbox}{HTML}{FCECEC}
\definecolor{greenbox}{HTML}{EDF8F0}
\definecolor{goldbox}{HTML}{FFF7E6}
\newtcolorbox{skillbox}[2][]{
    colback=#2, colframe=navy, boxrule=0.8pt, arc=2mm,
    left=2mm,right=2mm,top=1.5mm,bottom=1.5mm,
    fonttitle=\bfseries, title={#1}, halign title=center, breakable
}
\setlist[itemize]{leftmargin=1.2em,itemsep=0.35em,topsep=0.35em}
\setlist[enumerate]{leftmargin=1.4em,itemsep=0.35em,topsep=0.35em}
\newcolumntype{C}[1]{>{\centering\arraybackslash}p{#1}}
\newcolumntype{Y}{>{\centering\arraybackslash}X}

\begin{document}
\def\CourseName{AP Statistics}
\def\SchoolYear{2026--2027}
\def\MeetingLength{55 minutes}
\def\UnitNumberName{Unit 1: Exploring One-Variable Data \quad}
\def\LessonNumberName{Lesson 1.4: Representing a Categorical Variable with Graphs}

\begin{center}
    {\Large\bfseries \CourseName: \SchoolYear\ (\MeetingLength) \\
    {\small\bfseries \UnitNumberName \LessonNumberName}}
\end{center}

\begin{tcolorbox}[colback=sky,colframe=navy,boxrule=0.9pt,arc=2mm,
    left=2mm,right=2mm,top=1.5mm,bottom=1.5mm]
    \textbf{Primary Objective:} Students will construct and interpret bar charts and pie charts
    for categorical variables. Students will evaluate graphs for accuracy, appropriate labeling,
    and misleading features, and will explain when each graph type is most appropriate (Big Idea: UNC).
\end{tcolorbox}

\begin{skillbox}[Priority Ideas \& Skills]{goldbox}
\begin{minipage}[t]{0.48\linewidth}
    \textbf{AP Skill 2 --- Data Analysis}
    \begin{itemize}
        \item Construct a correctly labeled bar chart (with frequency or relative frequency on the vertical axis).
        \item Interpret bar charts and pie charts in context, describing what the display reveals about the distribution.
        \item Identify misleading features of a graph (e.g., non-zero axis, 3-D distortion, inconsistent scale).
    \end{itemize}
\end{minipage}
\hfill
\begin{minipage}[t]{0.48\linewidth}
    \textbf{Key Understandings}
    \begin{itemize}
        \item Bar charts are versatile and easy to compare; pie charts show part-to-whole relationships.
        \item The vertical axis of a bar chart should start at 0; otherwise the display is misleading.
        \item Unlike histograms (Lesson 1.5), bars in a bar chart do \textit{not} touch --- categories have no inherent order.
    \end{itemize}
\end{minipage}
\end{skillbox}

\begin{skillbox}[{Vocabulary, Concepts, \& Theorems}]{greenbox}
\begin{minipage}[t]{0.48\linewidth}
    \begin{itemize}
        \item \textbf{Bar chart (bar graph)} --- bars whose heights represent frequencies or relative frequencies for each category; bars do not touch
        \item \textbf{Pie chart} --- a circle divided into sectors proportional to relative frequencies
        \item \textbf{Segmented bar chart} --- bars divided into segments to show the composition of each group (useful for comparison)
    \end{itemize}
\end{minipage}
\hfill
\begin{minipage}[t]{0.48\linewidth}
    \begin{itemize}
        \item \textbf{Frequency bar chart} --- vertical axis shows counts
        \item \textbf{Relative frequency bar chart} --- vertical axis shows proportions or percentages (preferred for comparison)
        \item \textbf{Misleading graph} --- a visual that distorts or exaggerates differences through scale manipulation, 3-D effects, or inappropriate graph type
    \end{itemize}
\end{minipage}
\end{skillbox}

\begin{skillbox}[Activate Prior Knowledge \& Spiral Review]{sky}
\begin{minipage}[t]{0.48\linewidth}
    \textbf{Quick Review (3 min)}\\
    Display the relative frequency table from Lesson 1.3. Ask: ``If I want to \textit{see} the distribution quickly, what shape would I draw?'' Accept student suggestions. Today we formalize two of them.
\end{minipage}
\hfill
\begin{minipage}[t]{0.48\linewidth}
    \textbf{Spiral Connections}
    \begin{itemize}
        \item Relative frequencies from Lesson 1.3 become bar heights or pie-slice sizes.
        \item Segmented bar charts resurface in Unit 2 for comparing two categorical variables.
        \item AP exam free-response often asks students to construct \textit{and} interpret a graph --- both skills matter.
    \end{itemize}
\end{minipage}
\end{skillbox}

\begin{skillbox}[Hook \& Lesson]{sky}
\begin{minipage}[t]{0.48\linewidth}
    \textbf{Hook (5 min)}\\
    Display two ``bad'' graphs from real media (e.g., a bar chart with a non-zero y-axis, or a 3-D pie chart). Ask: ``What's wrong here? Does this graph make one option look more popular than it really is?'' Discuss that graphs can \textit{lie} even with correct data.
\end{minipage}
\hfill
\begin{minipage}[t]{0.48\linewidth}
    \textbf{Lesson Flow (15 min)}
    \begin{enumerate}
        \item Model constructing a bar chart by hand from a frequency table (title, labeled axes, correct scale starting at 0, bars not touching).
        \item Show both frequency and relative frequency versions of the same data.
        \item Demonstrate a pie chart: compute degree measures ($\text{relative frequency} \times 360^\circ$); label each sector.
        \item Compare: when is each graph type preferred?
    \end{enumerate}
\end{minipage}
\end{skillbox}

\begin{skillbox}[Explicit Instruction \& Active Monitoring]{sky}
\begin{minipage}[t]{0.48\linewidth}
    \textbf{Direct Instruction (10 min)}
    \begin{itemize}
        \item Emphasize the AP grading standard: a graph must have a title, labeled axes (with units), and a scale.
        \item Contrast bar chart vs.\ histogram: ``today's bars don't touch because categories are separate; next lesson, histogram bars touch because the data are continuous.''
        \item Walk through interpreting a bar chart: identify the modal category, notice large vs.\ small bars, and write a description.
    \end{itemize}
\end{minipage}
\hfill
\begin{minipage}[t]{0.48\linewidth}
    \textbf{Active Monitoring}
    \begin{itemize}
        \item As students sketch graphs in guided practice, check: Is the y-axis labeled? Does it start at 0? Are bars separated?
        \item Common error: making the bars touch (histogram instinct). Reinforce: categories $\to$ no touching.
        \item Ask: ``Is your graph truthful? Would it mislead a reader?''
    \end{itemize}
\end{minipage}
\end{skillbox}

\begin{skillbox}[Group Work \& Differentiation]{redbox}
\begin{minipage}[t]{0.48\linewidth}
    \textbf{Group Activity (8--10 min)}\\
    Groups receive a relative frequency table (preferred streaming service, 5 categories). Each group:
    \begin{enumerate}
        \item Constructs a relative frequency bar chart.
        \item Writes a 2-sentence interpretation.
        \item Identifies one change that would make the graph misleading.
        \item (Optional) Sketches a pie chart and computes sector angles.
    \end{enumerate}
\end{minipage}
\hfill
\begin{minipage}[t]{0.48\linewidth}
    \textbf{Differentiation}
    \begin{itemize}
        \item \textit{Support}: Provide a pre-drawn axis template; students only fill in bars and labels.
        \item \textit{Extension}: Students find and critique a real published graph (from news or social media) and rewrite the description as an AP-quality response.
        \item \textit{Technology}: If available, students create the graph in a spreadsheet app and compare to their hand-drawn version.
    \end{itemize}
\end{minipage}
\end{skillbox}

\begin{skillbox}[Individual Work \& Assessment]{redbox}
\begin{minipage}[t]{0.48\linewidth}
    \textbf{Exit Ticket (5 min)}\\
    A survey of 80 adults found: 32 watch news, 20 watch sports, 16 watch reality TV, 12 watch other. Students: (1) construct a relative frequency bar chart (sketch), (2) describe the distribution in one sentence, (3) identify one way someone could redraw the graph to be misleading.
\end{minipage}
\hfill
\begin{minipage}[t]{0.48\linewidth}
    \textbf{AP-Style Check}\\
    \textit{``Describe what the bar chart reveals about the preferred TV genre for these adults. Use the graph and relevant statistics in your response.''}\\[4pt]
    Assess: does the student \textit{refer to the graph}, include proportions, and use context?
\end{minipage}
\end{skillbox}

\begin{skillbox}[Reinforcement \& Extension]{goldbox}
\begin{minipage}[t]{0.48\linewidth}
    \textbf{Homework / Practice}
    \begin{itemize}
        \item Construct a bar chart from a provided frequency table and write a complete description.
        \item Critique one of the ``bad graph'' examples shown in class: explain what is misleading and how to fix it.
    \end{itemize}
\end{minipage}
\hfill
\begin{minipage}[t]{0.48\linewidth}
    \textbf{Extension / Preview}
    \begin{itemize}
        \item \textit{Extension}: Research ``lying with statistics'' or ``misleading graphs'' and bring one example to class.
        \item \textit{Preview}: Lesson 1.5 moves to \textit{quantitative} variables --- the displays look similar but have important differences (touching bars, meaningful order). Notice: what changes when data are numerical?
    \end{itemize}
\end{minipage}
\end{skillbox}

\end{document}
